\chapter{总结回顾}\label{chap:retrospective}

\section{项目组总结}

项目从四类角色的需求开始，验收时沿这四类角色走完同一笔订单。顾客选择外送或到店自取。外送单由商家接单、出餐，再交给骑手，配送异常由管理员处理。历史订单保留下单时的菜名和地址，模拟支付与退款有对应流水。系统接入地图导航和 AI 点餐，演示时可手工填写地址，也可用本地菜单筛选菜品。交付和验收结果见表~\ref{tab:final-metrics}。

\begin{longtable}{L{.258\linewidth}C{.255\linewidth}L{.448\linewidth}}
\caption{项目最终交付工程指标汇总表}\label{tab:final-metrics}\\
\toprule
\textbf{工程指标维度} & \textbf{实际交付量} & \textbf{可验证仓库依据} \\
\midrule
业务范围与用户故事 & 22 项（\mbox{113 故事点}） & 故事明细按 P0、P1 和 P2 登记。Sprint 计划共 118 点，记录完成率 100\%。 \\
前端视图与后端接口 & 33 视图 / \mbox{24 控制器} & 共 135 个 REST 端点与 4 套 CSS 变量主题 \\
数据库规模与迁移 & 23 张表 / \mbox{6 个迁移} & 含订单快照、资产账本、配送任务版本和地理坐标列 \\
版本协同与代码审查 & 265 提交 / \mbox{34 个 PR} & 覆盖 Sprint 1--3 与文档审计收尾阶段 \\
自动化与实机\mbox{质量验收} & 575 项测试\mbox{全通过} & 后端 340 项（53 类）+ 前端 235 项（28 文件）+ 12 缺陷闭环 \\
\bottomrule
\end{longtable}

排期以 22 项用户故事确定交付范围，一项故事通常涉及多个页面和接口。以配送订单为例，提交时先算金额，履约时再核对订单状态、任务归属和各角色看到的按钮。单元和集成测试验证规则，四端演示从顾客首次操作开始，确认订单能否走到送达。

\subsection{最有帮助的协作做法}

第一次四端联调时，页面都能打开，订单仍会卡在中间。要确认任务能否交给骑手，得回到商家出餐这一步，同时对照顾客看到的状态和骑手的待接列表。后来我们先沿同一笔订单确认状态变化、各页开放操作的条件，再回到各自代码。状态机上的一次转换，由三个角色在各自页面共同核对。

到店自取被起送价拦住时，我们在报错图旁记下履约方式、购物车金额、触发入口和预期结果。前端查结算页，后端查计价分支，测试再按同一条件复测。后来定位旧订单菜名、页面返回和地图缺少配置的问题，也用这套记法。修复人与复测人可以按台账在各自环境中重现问题。

回归和跨组验收把正常下单、接单之外的情形单列为用例，包括取消订单、数量超库存、登录失效和跨账号隔离。文档与实现有差异时，先核对角色、输入和保存结果，再按对应需求确定应补充的规格或代码改动。

\subsection{交付之后的改进顺序}

交付时发现的差异见附录表~\ref{tab:gap-improvement}。后续先统一需求、接口和页面上的订单状态名称，再复跑既有用例。高德和多模态服务接入有效密钥后，沿现有入口补测成功路径，并继续回归离线选址和本地推荐。并发下单和真实支付上线前，分别补测负载、支付回调和对账。已有用例、日志和配置可作为复测起点。

\section{项目组成员总结}

\subsection{范文丽（学号 3024244438）}

我负责前端开发、UI/UX 和数据看板，参与登录、首页、店铺详情、购物车、订单、AI 对话和偏好设置等页面，也维护商家工作台与管理员看板。我设计了清爽蓝白、深海夜色、薄荷青和暖阳橙四套 CSS 变量主题，并实现共享元素转场、响应式页面壳层和统计图表导出。

做登录页和首页动画时，我切到其他路由才发现全局样式也改了背景。屏幕缩到 375~px，登录框会溢出，关闭语音抽屉后音轨有时仍未释放。我们把样式限制在组件内，让动画遵循 \code{prefers-reduced-motion} 设置，再用 \code{theme.test.mjs} 和 \code{assistant-drawer.test.mjs} 复查。后来做看板，我先和后端确认统计口径与空数据展示，再画图。这一轮也让我熟悉了 Vue 3 的页面组织。

\subsection{李承文（学号 3024232129）}

我担任组长和 Scrum Master，负责后端、部署和跨端整合。我管理 GitHub 分支，跟进 34 个 Pull Request 的审查，并搭建 \code{docker-compose.demo.yml} 演示环境。开发任务包括 JWT 鉴权、订单提交事务、服务端计价、多资产结算、骑手配送、管理员治理、导航返回的 TDD 重构和部署解耦。

联调时，一笔取消单要退回资产，两位骑手抢同一单只能有一位成功，顾客在请求超时后重试也只能生成一张订单。我们针对这些情形补充状态前置条件、条件更新和事务，给每次改动留下可重复的测试。四人的课业时间很难排在一起，飞书清单和线上同步明确了阻塞事项和下一位接手者，PR 描述和测试记录也能查到上次改了什么。AI 帮忙写样板和查资料，资金和权限规则由组员共同确定。

\Needspace{7\baselineskip}
\subsection{赵紫怡（学号 3024244137）}

我担任产品负责人，负责需求分析、数据库设计和 AI 点餐开发。我维护 SRS 1.0 和 2.0，绘制系统上下文、参与者、用例和领域类模型，并整理飞书协同空间。我负责 23 张核心表和 \code{001}--\code{006} 增量迁移，参与历史商品名和地址快照设计，并实现 AI 会话、结构化推荐以及 DeepSeek 和 DashScope 的适配与降级。

SRS 的初稿先画正常流程。走查订单和助手时，我和组员补充了商家改名后的历史订单显示、查询他人订单时的助手响应，以及模型响应慢时数据库连接的释放。补上异常流程和状态后，大家可以对着同一操作讨论。多模态助手能把图或语音转成菜名，加购时还需核对菜单、库存和当前用户。和前端、测试一起验这些边界，是我印象最深的部分。

\subsection{仲禹潼（学号 3024244140）}

我负责测试和高德地图服务。内部走查使用 WPS 台账，9 月 16 日跨组验收记录了正常流程与边界场景，报告素材保留了其中 24 张截图。我协助完善前后端自动化回归，负责 POI 搜索、逆地理编码、GCJ-02 距离与送达时间估算、骑手导航，并验证缺少密钥时的处理。

刚接手测试时，我习惯打开页面看哪里有问题。自取卡在起送价、新人券可以重复领、游客点击返回又绕回原页，这几次让我发现，复现时必须写清账号角色和上一步操作。测地图也一样。坐标缺失时，从选址页查地址能否填写，从配送页查骑手能否读到文字地址。开发同学按记录复现，回归时也能据此分辨权限、数据和第三方服务的问题。

\paragraph{交付材料与复现}
归档时，团队将验收所用的源码版本 \code{6ddb219} 和运行环境记入 \code{verification.json}，将操作编号与截图一同保留。前端测试输出和后端测试摘要提供相应断言，台账记录已修缺陷的原因与提交。

复验可从顾客下单、商家接单、骑手接任务一直走到收货，记录各页状态。随后单独检查自取的费用和地址，以及历史明细中的菜名。把版本、账号角色和操作前后的状态一起归档，接手者就能按同样条件复测。这是我们希望继续保持的做法。需求写清楚，问题能复现，修复后留下一次还能运行的测试。
